\section{Methodology}

The scoping review will be conducted in accordance with the Joanna Briggs Institute (JBI) methodology for scoping reviews \citep{peters2020jbi}. The reporting will align with the Preferred Reporting Items for Systematic Reviews and Meta-Analyses extension for Scoping Reviews (PRISMA-ScR) guidelines \citep{tricco2018prisma}.

\subsection{Eligibility Criteria}
We will apply the Population (or participants), Concept, and Context (PCC) framework to define the eligibility criteria for our review. These criteria are summarized in Table~\ref{tab:pcc_criteria}.

\begin{table}[H]
\centering
\caption{PCC Eligibility Criteria Framework}
\label{tab:pcc_criteria}
\begin{tabularx}{\textwidth}{@{} p{2.5cm} L L @{}}
\toprule
\textbf{PCC Dimension} & \textbf{Inclusion Criteria} & \textbf{Exclusion Criteria} \\
\midrule
\textbf{Population} & Human patient cohorts of any age, gender, or clinical diagnosis eligible or evaluated for clinical trial inclusion. & Animal models, purely synthetic cohorts without reference clinical validation. \\
\addlinespace
\textbf{Concept} & Computerized or algorithmic systems utilizing AI (including machine learning, NLP, rule-based systems, and hybrid frameworks) specifically designed or evaluated for patient-to-trial matching. & Non-algorithmic recruitment strategies (e.g., manual screening, poster campaigns, digital advertising, simple keyword search without clinical context parsing). \\
\addlinespace
\textbf{Context} & All clinical environments (hospitals, multi-center trials, primary care, clinical research organizations) or secondary analysis of public clinical registries (e.g., ClinicalTrials.gov). & Non-clinical trial registries, general-purpose cohort building tools not evaluated for trial eligibility matching. \\
\addlinespace
\textbf{Study Types} & Peer-reviewed journal articles, conference papers (IEEE, ACM, etc.), and preprints (arXiv, bioRxiv) containing empirical evaluations. & Systematic reviews, letters, commentaries, editorials, and book reviews. Only English-language articles published from Jan 2018 to Aug 2026. \\
\bottomrule
\end{tabularx}
\end{table}

\subsection{Search Strategy}
A comprehensive search strategy will be executed across the following databases: MEDLINE (via PubMed), Embase, IEEE Xplore, and Scopus. The search strategy will be developed by the research team in collaboration with a medical librarian.

A sample search string designed for MEDLINE is shown in Box~\ref{box:search_strategy}. This search strategy was validated using test records from peer-reviewed literature in August 2026.

\begin{tcolorbox}[protocolbox, title={Box 2: Draft MEDLINE (PubMed) Search Strategy}, label=box:search_strategy]
\begin{minipage}{\textwidth}
\sffamily\small
\textbf{Search Query Block:}
\begin{tcolorbox}[searchbox]
\texttt{("Artificial Intelligence"[Mesh] OR "Natural Language Processing"[Mesh] OR "Machine Learning"[Mesh] OR "Deep Learning"[Mesh] OR "large language model*"[Title/Abstract] OR "NLP"[Title/Abstract] OR "neural network*"[Title/Abstract] OR "AI"[Title/Abstract])} \\{}
\texttt{AND} \\{}
\texttt{("Clinical Trials as Topic"[Mesh] OR "Patient Selection"[Mesh] OR "clinical trial*"[Title/Abstract] OR "trial eligibility"[Title/Abstract] OR "patient recruitment"[Title/Abstract])} \\{}
\texttt{AND} \\{}
\texttt{("match*"[Title/Abstract] OR "select*"[Title/Abstract] OR "screen*"[Title/Abstract] OR "cohort identification"[Title/Abstract] OR "automated parsing"[Title/Abstract])}
\end{tcolorbox}
\textbf{Filters:} English language, Publication date: January 1, 2018 -- August 31, 2026.
\end{minipage}
\end{tcolorbox}

\subsection{Study Selection and Screening}
Following the search, all identified citations will be uploaded to Covidence (or a similar screening software) and duplicates will be removed automatically. 

The screening process will consist of two distinct phases:
\begin{enumerate}
  \item \textbf{Title and Abstract Screening:} Two reviewers (SJ and AV) will independently screen the titles and abstracts of all retrieved studies against the predefined inclusion and exclusion criteria.
  \item \textbf{Full-Text Review:} Studies that pass the first phase will undergo full-text review by the same two reviewers independently.
\end{enumerate}
Any discrepancies or disagreements between the reviewers at either stage will be discussed and resolved through consensus. If consensus cannot be reached, a third reviewer (CL) will make the final decision. The results of the search and study screening process will be reported in full in the final scoping review and presented using a PRISMA 2020 flow diagram.
